\section{Preparación de las VPS}\label{sec:preparacion}

\subsection{Primer acceso por SSH}
\pendiente{conexión, cambio de contraseña root, recomendación de llaves SSH}

\subsection{Actualización del sistema y zona horaria}
\begin{lstlisting}
# apt update && apt upgrade -y
# timedatectl set-timezone America/Mexico_City
\end{lstlisting}

\subsection{Nombre de host}
\pendiente{vmi3633025 (BD), vmi3633021 (Backend), vmi3633026 (Frontend)}

\subsection{Instalación de Podman}
\pendiente{\texttt{apt install podman}, verificación con \texttt{podman info}, registro por defecto docker.io}

\subsection{Habilitar el reinicio automático de contenedores}
\pendiente{\texttt{podman-restart.service} y la política \texttt{--restart always}}

\subsection{Estructura de directorios en el servidor}
\pendiente{carpeta de trabajo de cada VPS (scripts, \texttt{config-files/}, \texttt{jars/}, \texttt{database-scripts/}, respaldos)}

\subsection{Transferencia de archivos al servidor}
\pendiente{uso de \texttt{scp} / WinSCP / FileZilla desde la carpeta \texttt{Produccion}}

\subsection{Consideraciones de firewall}
\pendiente{por qué ufw no basta con Podman (el DNAT lo omite) y por qué se usan reglas en la tabla raw de iptables limitadas a la interfaz pública}
